\section{Normalization of Expression Values}
    Raw gene expression data is normalized in the following steps:
    \begin{enumerate}
      \item Division by gene-wise standard deviation for scale normalization. For estimation of the correct standard deviation in low dimensional data-sets see below (\ref{loess_sd_estimation}).
        \item Quantile normalization across samples to reduce global expression biases.
        \item Log-transformation: \( x \mapsto \log(1 + x) \) to stabilize variance and suppress outliers.
        \item For stress response studies, log fold-changes are computed: 
        \[
            \Delta_{\text{stress}} = \log(1 + x_{\text{stress}}) - \log(1 + x_{\text{control}})
        \]
        This effectively makes each stress related axis show the fold change in
        gene expression under stress, e.g.\ ``drought in root divided by control
        root''.
    \end{enumerate}

    \subsection{Variance (Standard Deviation) Stabilization for Scaling}\label{loess_sd_estimation}

        In Tensor Omics, normalization of expression vectors requires division by gene-wise standard deviation (SD) across samples.
        However, empirical SD estimates are unstable for small sample sizes.
        To address this, we implement a dimension-dependent stabilization strategy that improves robustness while maintaining geometric interpretability.

        The strategy is defined as follows:

        \begin{itemize}
            \item For small sample sizes \( n < 10 \):
            \begin{itemize}
                \item Empirical SD estimates are highly unreliable.
                \item Therefore, for each gene, the standard deviation is replaced by a
                  LOESS-fitted value (see~\ref{gene_sd_loess_alg}).
                \item The LOESS curve is fitted to model the relationship between
                  gene-wise mean expression and gene-wise empirical SD across all
                  genes.
            \end{itemize}
            
            \item For intermediate sample sizes \( 10 \leq n < 21 \):
            \begin{itemize}
                \item A weighted combination of empirical SD and LOESS-fitted SD is used.
                \item The effective standard deviation is computed as
                      \[
                          \text{SD}_{\text{effective}} = \lambda \times \text{SD}_{\text{empirical}} + (1-\lambda) \times \text{SD}_{\text{LOESS}},
                      \]
                where \( \lambda = (n-10)/11 \) increases linearly from 0 at \( n=10 \) to 1 at \( n=21 \).
            \end{itemize}
            
            \item For large sample sizes \( n \geq 21 \):
            \begin{itemize}
                \item Empirical SD estimates are considered sufficiently stable.
                \item Pure empirical SD is used without adjustment.
            \end{itemize}
        \end{itemize}

        This stabilization ensures that scaling by SD is consistent and resistant to noise across a wide range of dataset sizes.
        It preserves the empirical geometric structure of the data while suppressing variance artifacts that would otherwise distort Tensor Omics analyses, particularly in small or medium sample regimes.

        \subsubsection{LOESS Estimation of Gene-wise Standard Deviation}\label{gene_sd_loess_alg}

        To obtain stabilized estimates of gene-wise standard deviations for small
        sample sizes, Tensor Omics applies a LOESS smoothing procedure based on the
        relationship between gene-wise mean expression and gene-wise empirical SD.

        The procedure consists of the following steps:

        \begin{itemize}
            \item Calculate the empirical mean and empirical standard deviation for
              each gene across all samples.
            \item Plot each gene as a point in a scatter plot with
            \begin{itemize}
                \item X-axis: gene-wise mean expression,
                \item Y-axis: gene-wise empirical standard deviation.
            \end{itemize}
            \item Fit a LOESS curve to these points to model the typical dependence of
              standard deviation on mean expression.
            \item For each gene, given its mean expression value, predict its standard
              deviation by evaluating the fitted LOESS curve at the gene's mean.
        \end{itemize}

        This method exploits the fact that gene-wise mean expression values are
        substantially more stable than standard deviation estimates, especially at low
        sample sizes. While the empirical standard deviation suffers from high
        variability when few data points are available, the mean value remains a more
        robust and lower-variance statistic. By basing the standard deviation
        estimation on the more reliable mean, the procedure suppresses noise-induced
        artifacts and yields smoothed, biologically plausible variance estimates. This
        LOESS-based adjustment is critical for preserving the integrity of the
        normalization process in small-sample regimes, ensuring that subsequent
        geometric analyses remain meaningful and reproducible.